
\subsubsection{3.1 Confidence Operationalization}

We operationalize model confidence using the logit margin:

$$\text{margin} = \text{logit}_{(1)} - \text{logit}_{(2)}$$

where $\text{logit}_{(1)}$ and $\text{logit}_{(2)}$ are the highest and second-highest logits among answer options. This operationalization is standard for multiple-choice tasks where probability mass concentrates on answer tokens.

\subsubsection{3.2 Discriminative Degradation (H-E1)}

We compute AUROC for the binary classification task of predicting correctness from margin:

$$\text{AUROC} = P(\text{margin}_{\text{correct}} > \text{margin}_{\text{incorrect}})$$

This measures how well margin values rank predictions by correctness probability, independent of calibration. We compute bootstrap confidence intervals (n=1,000) for AUROC differences between base and instruction-tuned models.

\subsubsection{3.3 Conditional Margin Analysis (H-M1)}

We analyze conditional margin distributions separately for correct and incorrect predictions:

$$E[\text{margin} | \text{incorrect}]_{\text{instruct}} \quad \text{vs.} \quad E[\text{margin} | \text{incorrect}]_{\text{base}}$$

If instruction tuning inflates margins disproportionately for incorrect predictions, this would explain AUROC degradation. We use permutation tests (n=9,999) to assess statistical significance and compute Cohen's d for effect size.

\subsubsection{3.4 Percentile-Normalized Monotonicity (H-M2)}

To separate scale effects from shape effects, we apply percentile normalization:

\begin{enumerate}
\item Transform margins to percentile ranks within each model
\item Convert to z-scores using the inverse normal CDF
\item Fit logistic regression: $P(\text{correct}) = \sigma(\alpha + \beta \cdot z(\text{margin}))$
\item Compare $\beta_{\text{base}}$ vs. $\beta_{\text{instruct}}$
\end{enumerate}

If $\beta_{\text{instruct}} < \beta_{\text{base}}$ after normalization, the confidence-correctness relationship is attenuated in shape, not merely scale. We use paired bootstrap (n=1,000) to compute confidence intervals for the slope difference.

\subsubsection{3.5 Brier Score Decomposition (H-M3)}

We apply Murphy's (1973) decomposition:

$$\text{Brier} = \text{Reliability} - \text{Refinement} + \text{Uncertainty}$$

where Reliability measures calibration error, Refinement measures discrimination, and Uncertainty measures base rate entropy. If RLHF causes geometric distortion (affecting shape), Refinement should decrease. If RLHF causes scalar distortion (temperature-like), Reliability should change while Refinement remains stable.
